% Options for packages loaded elsewhere
\PassOptionsToPackage{unicode}{hyperref}
\PassOptionsToPackage{hyphens}{url}
\documentclass[
]{article}
\usepackage{xcolor}
\usepackage{amsmath,amssymb}
\setcounter{secnumdepth}{-\maxdimen} % remove section numbering
\usepackage{iftex}
\ifPDFTeX
  \usepackage[T1]{fontenc}
  \usepackage[utf8]{inputenc}
  \usepackage{textcomp} % provide euro and other symbols
\else % if luatex or xetex
  \usepackage{unicode-math} % this also loads fontspec
  \defaultfontfeatures{Scale=MatchLowercase}
  \defaultfontfeatures[\rmfamily]{Ligatures=TeX,Scale=1}
\fi
\usepackage{lmodern}
\ifPDFTeX\else
  % xetex/luatex font selection
\fi
% Use upquote if available, for straight quotes in verbatim environments
\IfFileExists{upquote.sty}{\usepackage{upquote}}{}
\IfFileExists{microtype.sty}{% use microtype if available
  \usepackage[]{microtype}
  \UseMicrotypeSet[protrusion]{basicmath} % disable protrusion for tt fonts
}{}
\makeatletter
\@ifundefined{KOMAClassName}{% if non-KOMA class
  \IfFileExists{parskip.sty}{%
    \usepackage{parskip}
  }{% else
    \setlength{\parindent}{0pt}
    \setlength{\parskip}{6pt plus 2pt minus 1pt}}
}{% if KOMA class
  \KOMAoptions{parskip=half}}
\makeatother
\setlength{\emergencystretch}{3em} % prevent overfull lines
\providecommand{\tightlist}{%
  \setlength{\itemsep}{0pt}\setlength{\parskip}{0pt}}
\usepackage[]{natbib}
\bibliographystyle{plainnat}
\usepackage{bookmark}
\IfFileExists{xurl.sty}{\usepackage{xurl}}{} % add URL line breaks if available
\urlstyle{same}
\hypersetup{
  hidelinks,
  pdfcreator={LaTeX via pandoc}}

\author{}
\date{}

\begin{document}

Results included in this manuscript come from preprocessing performed
using \emph{fMRIPrep} 25.2.5 (\citet{fmriprep1}; \citet{fmriprep2};
RRID:SCR\_016216), which is based on \emph{Nipype} 1.10.0
(\citet{nipype1}; \citet{nipype2}; RRID:SCR\_002502).

\begin{description}
\item[Preprocessing of B0 inhomogeneity mappings]
A total of 1 fieldmaps were found available within the input BIDS
structure for this particular subject. A \emph{B0} nonuniformity map (or
\emph{fieldmap}) was estimated from the phase-drift map(s) measure with
two consecutive GRE (gradient-recalled echo) acquisitions. The
corresponding phase-map(s) were phase-unwrapped with \texttt{prelude}
(FSL None).
\item[Anatomical data preprocessing]
A total of 1 T1-weighted (T1w) images were found within the input BIDS
dataset. The T1w image was corrected for intensity non-uniformity (INU)
with \texttt{N4BiasFieldCorrection} \citep{n4}, distributed with ANTs
2.6.2 \citep[RRID:SCR\_004757]{ants}, and used as T1w-reference
throughout the workflow. The T1w-reference was then skull-stripped with
a \emph{Nipype} implementation of the \texttt{antsBrainExtraction.sh}
workflow (from ANTs), using OASIS30ANTs as target template. Brain tissue
segmentation of cerebrospinal fluid (CSF), white-matter (WM) and
gray-matter (GM) was performed on the brain-extracted T1w using
\texttt{fast} \citep[FSL (version unknown),
RRID:SCR\_002823,][]{fsl_fast}. Volume-based spatial normalization to
two standard spaces (MNI152NLin6Asym, MNI152NLin2009cAsym) was performed
through nonlinear registration with \texttt{antsRegistration} (ANTs
2.6.2), using brain-extracted versions of both T1w reference and the T1w
template. The following templates were were selected for spatial
normalization and accessed with \emph{TemplateFlow}
\citep[25.0.4,][]{templateflow}: \emph{FSL's MNI ICBM 152 non-linear 6th
Generation Asymmetric Average Brain Stereotaxic Registration Model}
{[}\citet{mni152nlin6asym}, RRID:SCR\_002823; TemplateFlow ID:
MNI152NLin6Asym{]}, \emph{ICBM 152 Nonlinear Asymmetrical template
version 2009c} {[}\citet{mni152nlin2009casym}, RRID:SCR\_008796;
TemplateFlow ID: MNI152NLin2009cAsym{]}.
\item[Functional data preprocessing]
For each of the 7 BOLD runs found per subject (across all tasks and
sessions), the following preprocessing was performed. First, a reference
volume was generated from the shortest echo of the BOLD run, using a
custom methodology of \emph{fMRIPrep}, for use in head motion
correction. Head-motion parameters with respect to the BOLD reference
(transformation matrices, and six corresponding rotation and translation
parameters) are estimated before any spatiotemporal filtering using
\texttt{mcflirt} \citep[FSL ,][]{mcflirt}. The estimated \emph{fieldmap}
was then aligned with rigid-registration to the target EPI (echo-planar
imaging) reference run. The field coefficients were mapped on to the
reference EPI using the transform. The BOLD reference was then
co-registered to the T1w reference using \texttt{mri\_coreg}
(FreeSurfer) followed by \texttt{flirt} \citep[FSL ,][]{flirt} with the
boundary-based registration \citep{bbr} cost-function. Co-registration
was configured with six degrees of freedom. BOLD runs were slice-time
corrected to 0.75s (0.5 of slice acquisition range 0s-1.5s) using
\texttt{3dTshift} from AFNI \citep[RRID:SCR\_005927]{afni}. A T2★ map
was estimated from the preprocessed EPI echoes using tedana's t2smap
workflow \citep{DuPre2021}, by voxel-wise fitting the maximal number of
echoes with reliable signal in that voxel to a monoexponential signal
decay model with nonlinear regression. The T2★/S0 estimates from a
log-linear regression fit were used for initial values. The calculated
T2★ map was then used to optimally combine preprocessed BOLD across
echoes following the method described in \citep{posse_t2s}. The
optimally combined time series was carried forward as the
\emph{preprocessed BOLD}. Several confounding time-series were
calculated based on the \emph{preprocessed BOLD}: framewise displacement
(FD), DVARS and three region-wise global signals. FD was computed using
two formulations following Power (absolute sum of relative motions,
\citet{power_fd_dvars}) and Jenkinson (relative root mean square
displacement between affines, \citet{mcflirt}). FD and DVARS are
calculated for each functional run, both using their implementations in
\emph{Nipype} \citep[following the definitions by][]{power_fd_dvars}.
The three global signals are extracted within the CSF, the WM, and the
whole-brain masks. Additionally, a set of physiological regressors were
extracted to allow for component-based noise correction
\citep[\emph{CompCor},][]{compcor}. Principal components are estimated
after high-pass filtering the \emph{preprocessed BOLD} time-series
(using a discrete cosine filter with 128s cut-off) for the two
\emph{CompCor} variants: temporal (tCompCor) and anatomical (aCompCor).
tCompCor components are then calculated from the top 2\% variable voxels
within the brain mask. For aCompCor, three probabilistic masks (CSF, WM
and combined CSF+WM) are generated in anatomical space. The
implementation differs from that of Behzadi et al.~in that instead of
eroding the masks by 2 pixels on BOLD space, a mask of pixels that
likely contain a volume fraction of GM is subtracted from the aCompCor
masks. This mask is obtained by thresholding the corresponding partial
volume map at 0.05, and it ensures components are not extracted from
voxels containing a minimal fraction of GM. Finally, these masks are
resampled into BOLD space and binarized by thresholding at 0.99 (as in
the original implementation). Components are also calculated separately
within the WM and CSF masks. For each CompCor decomposition, the
\emph{k} components with the largest singular values are retained, such
that the retained components' time series are sufficient to explain 50
percent of variance across the nuisance mask (CSF, WM, combined, or
temporal). The remaining components are dropped from consideration. The
head-motion estimates calculated in the correction step were also placed
within the corresponding confounds file. The confound time series
derived from head motion estimates and global signals were expanded with
the inclusion of temporal derivatives and quadratic terms for each
\citep{confounds_satterthwaite_2013}. Frames that exceeded a threshold
of 0.5 mm FD or 1.5 standardized DVARS were annotated as motion
outliers. Additional nuisance timeseries are calculated by means of
principal components analysis of the signal found within a thin band
(\emph{crown}) of voxels around the edge of the brain, as proposed by
\citep{patriat_improved_2017}. All resamplings can be performed with
\emph{a single interpolation step} by composing all the pertinent
transformations (i.e.~head-motion transform matrices, susceptibility
distortion correction when available, and co-registrations to anatomical
and output spaces). Gridded (volumetric) resamplings were performed
using \texttt{nitransforms}, configured with cubic B-spline
interpolation.
\item[Functional data preprocessing]
For each of the 7 BOLD runs found per subject (across all tasks and
sessions), the following preprocessing was performed. First, a reference
volume was generated from the shortest echo of the BOLD run, using a
custom methodology of \emph{fMRIPrep}, for use in head motion
correction. Head-motion parameters with respect to the BOLD reference
(transformation matrices, and six corresponding rotation and translation
parameters) are estimated before any spatiotemporal filtering using
\texttt{mcflirt} \citep[FSL ,][]{mcflirt}. The estimated \emph{fieldmap}
was then aligned with rigid-registration to the target EPI (echo-planar
imaging) reference run. The field coefficients were mapped on to the
reference EPI using the transform. The BOLD reference was then
co-registered to the T1w reference using \texttt{mri\_coreg}
(FreeSurfer) followed by \texttt{flirt} \citep[FSL ,][]{flirt} with the
boundary-based registration \citep{bbr} cost-function. Co-registration
was configured with six degrees of freedom. BOLD runs were slice-time
corrected to 0.751s (0.5 of slice acquisition range 0s-1.5s) using
\texttt{3dTshift} from AFNI \citep[RRID:SCR\_005927]{afni}. A T2★ map
was estimated from the preprocessed EPI echoes using tedana's t2smap
workflow \citep{DuPre2021}, by voxel-wise fitting the maximal number of
echoes with reliable signal in that voxel to a monoexponential signal
decay model with nonlinear regression. The T2★/S0 estimates from a
log-linear regression fit were used for initial values. The calculated
T2★ map was then used to optimally combine preprocessed BOLD across
echoes following the method described in \citep{posse_t2s}. The
optimally combined time series was carried forward as the
\emph{preprocessed BOLD}. Several confounding time-series were
calculated based on the \emph{preprocessed BOLD}: framewise displacement
(FD), DVARS and three region-wise global signals. FD was computed using
two formulations following Power (absolute sum of relative motions,
\citet{power_fd_dvars}) and Jenkinson (relative root mean square
displacement between affines, \citet{mcflirt}). FD and DVARS are
calculated for each functional run, both using their implementations in
\emph{Nipype} \citep[following the definitions by][]{power_fd_dvars}.
The three global signals are extracted within the CSF, the WM, and the
whole-brain masks. Additionally, a set of physiological regressors were
extracted to allow for component-based noise correction
\citep[\emph{CompCor},][]{compcor}. Principal components are estimated
after high-pass filtering the \emph{preprocessed BOLD} time-series
(using a discrete cosine filter with 128s cut-off) for the two
\emph{CompCor} variants: temporal (tCompCor) and anatomical (aCompCor).
tCompCor components are then calculated from the top 2\% variable voxels
within the brain mask. For aCompCor, three probabilistic masks (CSF, WM
and combined CSF+WM) are generated in anatomical space. The
implementation differs from that of Behzadi et al.~in that instead of
eroding the masks by 2 pixels on BOLD space, a mask of pixels that
likely contain a volume fraction of GM is subtracted from the aCompCor
masks. This mask is obtained by thresholding the corresponding partial
volume map at 0.05, and it ensures components are not extracted from
voxels containing a minimal fraction of GM. Finally, these masks are
resampled into BOLD space and binarized by thresholding at 0.99 (as in
the original implementation). Components are also calculated separately
within the WM and CSF masks. For each CompCor decomposition, the
\emph{k} components with the largest singular values are retained, such
that the retained components' time series are sufficient to explain 50
percent of variance across the nuisance mask (CSF, WM, combined, or
temporal). The remaining components are dropped from consideration. The
head-motion estimates calculated in the correction step were also placed
within the corresponding confounds file. The confound time series
derived from head motion estimates and global signals were expanded with
the inclusion of temporal derivatives and quadratic terms for each
\citep{confounds_satterthwaite_2013}. Frames that exceeded a threshold
of 0.5 mm FD or 1.5 standardized DVARS were annotated as motion
outliers. Additional nuisance timeseries are calculated by means of
principal components analysis of the signal found within a thin band
(\emph{crown}) of voxels around the edge of the brain, as proposed by
\citep{patriat_improved_2017}. All resamplings can be performed with
\emph{a single interpolation step} by composing all the pertinent
transformations (i.e.~head-motion transform matrices, susceptibility
distortion correction when available, and co-registrations to anatomical
and output spaces). Gridded (volumetric) resamplings were performed
using \texttt{nitransforms}, configured with cubic B-spline
interpolation.
\end{description}

Many internal operations of \emph{fMRIPrep} use \emph{Nilearn} 0.11.1
\citep[RRID:SCR\_001362]{nilearn}, mostly within the functional
processing workflow. For more details of the pipeline, see
\href{https://fmriprep.readthedocs.io/en/latest/workflows.html}{the
section corresponding to workflows in \emph{fMRIPrep}'s documentation}.

\subsubsection{Copyright Waiver}\label{copyright-waiver}

The above boilerplate text was automatically generated by fMRIPrep with
the express intention that users should copy and paste this text into
their manuscripts \emph{unchanged}. It is released under the
\href{https://creativecommons.org/publicdomain/zero/1.0/}{CC0} license.

\subsubsection{References}\label{references}

\bibliography{/out/logs/CITATION.bib}

\end{document}
